\subsection{Perimeter-budget optimization}
\label{subsec:app-perimeter-budget-optimization}

\paragraph{Exact interface.}
The center and V-triangle boundary contributions are those of
Theorem~\ref{thm:boundary-trace-table}.  The public routing consequence is
Proposition~\ref{prop:length-branches}.

\paragraph{Variables and feasible set.}
The signed center variables range over the CE1 or CE2 cells of
\zcref{prop:signed-center-normal-form}.  The remaining variables
are the numbers and strict caps of the classified V triangles in the stated
routing branch.

\paragraph{Objective.}
Maximize the total perimeter trace subject to those caps and compare it with
$\mathcal H^1(\partial H)=6$.

\paragraph{Exact result.}
The exact deficit, small-slack, endpoint-slack, and routing conclusions are
the statements of the four results below.

\paragraph{Solution.}
Their proofs perform the signed substitutions and strict endpoint checks.

\begin{lemma}[Master perimeter deficit]
\label{lem:signed-perimeter-deficit}
Put
\[
 \theta=\frac2{\sqrt3}-1.
\]
Suppose a branch has a center contribution $L_C$, $n$ supercritical Vd0
V triangles, $m\ge1$ distinguished nonsupercritical V triangles with strict boundary caps
$q_1,\ldots,q_m<1$, and $6-n-m$ remaining V triangles whose boundary contributions
are at most $1$. If
\begin{equation}
 L_C+n\theta\le\sum_{j=1}^m(1-q_j),
 \label{eq:signed-perimeter-deficit}
\end{equation}
then the seven triangles do not cover $\partial H$.
\end{lemma}

\begin{proof}
The supercritical Vd0 cap is $2/\sqrt3$. Hence the total boundary contribution
is strictly less than
\[
 \begin{aligned}
 L_C+n\frac2{\sqrt3}+\sum_{j=1}^m q_j+(6-n-m)
 &=6+L_C+n\theta-\sum_{j=1}^m(1-q_j)\\
 &\le6.
 \end{aligned}
\]
The strictness follows from $m\ge1$ and the strict distinguished-V triangle caps. A
cover of the perimeter would give the opposite weak inequality by
subadditivity.
\end{proof}

\begin{lemma}[Small center slacks in the one-Vd1/Vd2 branch]
\label{lem:signed-small-slack}
Assume there is exactly one supercritical Vd0 V triangle, exactly one Vd1/Vd2 V triangle,
and four nonsupercritical Vd0 V triangles. Then every candidate surviving the
perimeter budget is CE2 and satisfies
\begin{equation}
 \boxed{
 \alpha+\delta<\frac1{24},
 \qquad
 \alpha+\delta<\frac{\min\{R,W\}}6.}
 \label{eq:signed-small-slacks}
\end{equation}
\end{lemma}

\begin{proof}
The Vd1/Vd2 boundary cap is strict and equals $1/2$. If the perimeter were
covered, Lemma~\ref{lem:signed-perimeter-deficit} would force
\[
 L_C+\theta>\frac12.
\]
Put
\[
 \kappa=\frac12-\theta=\frac32-\frac2{\sqrt3}.
\]
Thus $L_C>\kappa$. The CE1 cap in
\zcref{cor:signed-center-boundary} gives
\[
 L_C+\theta
 \le\left(\frac{\sqrt3}{2}-\frac34\right)
   +\left(\frac2{\sqrt3}-1\right)
 <\frac12,
\]
so every survivor is CE2.

Set $T=\alpha+\delta$. The CE2 formula gives
\[
 T<M(E):=
 \frac{RW}{E^2}\left(\frac{E}{1+E}-\kappa\right).
\]
Since $RW=1-E^2$,
\[
 M(E)=\frac{1-E}{E^2}\bigl(E-\kappa(1+E)\bigr),
\]
and
\[
 M'(E)=-\frac{E-2\kappa}{E^3}<0
 \qquad\left(\frac{\sqrt3}{2}\le E<1\right).
\]
Here $2\kappa<\sqrt3/2$, because it is equivalent to
$6\sqrt3<11$, whose square is $108<121$. Therefore
\[
 T<M\left(\frac{\sqrt3}{2}\right)
 =\frac{8\sqrt3}{9}-\frac32<\frac1{24}.
\]
The last inequality is equivalent to $64\sqrt3<111$, and follows after
squaring from $12288<12321$.

For the second estimate, use $E/(1+E)<1/2$:
\[
 T<\frac{RW}{E^2}\left(\frac12-\kappa\right)
   =\frac{RW}{E^2}\theta.
\]
Since
\[
 E^2=W+R^2=R+W^2,
\]
one has $RW/E^2\le R$ and $RW/E^2\le W$. Also
$\theta<1/6$, equivalently $12<7\sqrt3$, whose square is $144<147$.
This proves \eqref{eq:signed-small-slacks}.
\end{proof}

\begin{lemma}[CE2 initial endpoints dominate the total slack]
\label{lem:signed-endpoints-dominate-slack}
Assume $\Delta_R>0$ and $\Delta_L>0$. Put
\[
 T=\alpha+\delta,
 \qquad
 x=\frac{\eta+T}{W},
 \qquad
 y=\frac{\eta+T}{R}.
\]
Then
\begin{equation}
 \boxed{T<\frac12\min\{x,y\}.}
 \label{eq:signed-endpoint-slack}
\end{equation}
\end{lemma}

\begin{proof}
\zcref{eq:signed-total-slack} gives $T<\eta/E$. Moreover
\[
 E^2-(2R-1)^2=3RW>0,
\]
so $|2R-1|<E$ and therefore $|2R-1|T<\eta$. Hence
\[
 x-2T=\frac{\eta-(1-2R)T}{W}>0,
\]
and
\[
 y-2T=\frac{\eta-(2R-1)T}{R}>0.
\]
This is \eqref{eq:signed-endpoint-slack}.
\end{proof}

\begin{corollary}[Length branches closed by the common budgets]
\label{cor:signed-budget-branches}
The following branches are impossible.
\begin{enumerate}
\item $N_{\mathrm{gap}}=0$, $N_+=1$, and at least one Vd1/Vd2 V triangle;
\item $N_+=0$ and at least one Vd1/Vd2 V triangle, for either CE1 or CE2;
\item CE1, $N_+=1$, and at least one Vd1/Vd2 V triangle;
\item CE2, $N_+=1$, and at least two Vd1/Vd2 V triangles;
\item CE2, $N_+=1$, with a Vd2 triangle $T_i$ containing $M_{i-1}$ or
$M_{i+1}$;
\item CE1 or CE2 with $N_++N_{\mathrm{sp}}\ge3$;
\item CE1 or CE2 with $N_+\ge2$.
\end{enumerate}
\end{corollary}

\begin{proof}
For item 1, Lemma~\ref{lem:gap-exhaustion} says that the six open V
triangles cover $\partial H$.  The unique supercritical triangle is Vd0, so
the supercritical cap, one strict Vd1/Vd2 cap, and four unit caps give
\[
 \frac2{\sqrt3}+\frac12+4<6.
\]

For item 2, the center contribution is strictly less than $1/2$ and the
chosen Vd1/Vd2 V triangle has strict cap $1/2$; the other five V triangles contribute at
most $1$ each. This is Lemma~\ref{lem:signed-perimeter-deficit} with
$n=0$ and $q_1=1/2$.

For item 3, the unique supercritical V triangle is Vd0, while the CE1 center cap,
the Vd1/Vd2 cap, and the four remaining unit caps give
\[
 \left(\frac{\sqrt3}{2}-\frac34\right)
 +\frac12+\frac2{\sqrt3}+4<6.
\]
For item 4, the corresponding CE2 sum is strictly less than
\[
 \frac12+\frac12+\frac12+\frac2{\sqrt3}+3<6.
\]
For item 5, the $M_{i-1}$/$M_{i+1}$ Vd2 cap is $1/3$, and the total is
strictly less than
\[
 \frac12+\frac13+\frac2{\sqrt3}+4<6.
\]

Item 6 is Theorem~\ref{thm:common-skeleton-count}.

For item 7, Lemma~\ref{lem:positive-support-rescuer} supplies a
positive-support triangle distinct from two supercritical triangles.  Hence
$N_++N_{\mathrm{sp}}\ge3$, and
Theorem~\ref{thm:common-skeleton-count} applies.
\end{proof}

\begin{lemma}[Exact trace-routing arithmetic]
\label{lem:app-length-branch-calculation}
Every configuration satisfying one of the six hypotheses in
Proposition~\ref{prop:length-branches} has total available trace strictly
smaller than the target trace required in that proposition.
\end{lemma}

\begin{proof}
For item~(1), Lemma~\ref{lem:gap-exhaustion} says that the six open V
triangles cover $\partial H$.  Their incident open traces overlap on every
edge, so
\[
 1<B_i+A_{i+1}\qquad(0\le i\le5).
\]
Summing gives $6<\sum_i(A_i+B_i)$, whereas $N_+=0$ gives the opposite weak
inequality.

Item~(2) is item~(1) of Corollary~\ref{cor:signed-budget-branches}.

Item~(3) is Theorem~\ref{thm:common-skeleton-count}.  Items~(4) and~(5) are
items~(2) and~(3), respectively, of
Corollary~\ref{cor:signed-budget-branches}.  Item~(6) is item~(5) of that
corollary, using Lemma~\ref{lem:vd2-neighbor-midpoint-cap}.
\end{proof}

\paragraph{Body consequence.}
The preceding solved budget problem supplies
Proposition~\ref{prop:length-branches}, including its center-independent
$N_{\mathrm{gap}}=0$ rows.
